\begin{exercise}[Age and lookback time in Einstein--de Sitter]{Tong Cosmology, Cambridge Part III Sheet 1}
For a flat matter-only universe, derive \(t_0\), \(t(z)\), and the lookback
time as functions of \(H_0\) and \(z\).  Compare the result with the Hubble
time.
\end{exercise}

\begin{exercise}[Where angular size stops decreasing]{Cambridge Part III Cosmology, examples}
Find the redshift at which the angular-diameter distance is maximal in an
Einstein--de Sitter universe.  Explain why a standard ruler begins to appear
larger beyond this redshift.
\end{exercise}

\begin{exercise}[The Alcock--Paczynski test]{Cambridge Part III Cosmology, 2022 exam}
A statistically spherical comoving structure is observed with angular extent
\(\Delta\theta\) and redshift extent \(\Delta z\).  Derive the
Alcock--Paczynski combination \(F_{\rm AP}(z)\) that should make its radial and
transverse sizes agree.
\end{exercise}

\begin{exercise}[Cosmography through the jerk]{Cambridge Part III Cosmology, 2018 exam}
Extend the low-redshift luminosity-distance expansion through \(z^3\), using
the deceleration \(q_0\), jerk \(j_0\), and curvature.  Identify which
combination first distinguishes flat \(\Lambda\)CDM from a general smooth
expansion history.
\end{exercise}

\begin{exercise}[Weak-lensing convergence in Robertson--Walker spacetime]{Cambridge Part III Cosmology, examples}
In the Born approximation, derive the convergence produced by a weak density
contrast along the line of sight to a source at comoving distance \(\chi_s\).
Identify the geometric lensing kernel.
\end{exercise}

\begin{exercise}[Source counts per redshift]{Tong Cosmology, Cambridge Part III Sheet 2}
For comoving number density \(n_c(z)\) and selection probability \(S(z)\),
derive \(\dd N/(\dd z\,\dd\Omega)\).  Explain separately the roles of geometry,
population evolution, and selection.
\end{exercise}

\begin{exercise}[\(K\)-correction for a power-law spectrum]{Cambridge Part III Cosmology, examples}
A source has spectral luminosity \(L_\nu\propto\nu^{-\alpha}\).  Relate the
observed flux density at fixed observing frequency to \(L_\nu\) and
luminosity distance, and derive the corresponding power-law \(K\)-correction.
\end{exercise}

\begin{exercise}[Cosmic chronometers]{Cambridge Part III Cosmology, examples}
Show that differential ages of passively evolving galaxies give
\[
H(z)=-\frac{1}{1+z}\frac{\dd z}{\dd t}.
\]
List the assumptions needed for this to be a geometric expansion-rate
measurement.
\end{exercise}

\begin{exercise}[Radial and transverse baryon acoustic scales]{Tong Cosmology, Cambridge Part III Sheet 4}
A comoving standard ruler has length \(r_s\).  Derive its observed angular
size and redshift separation at redshift \(z\).  Show which combinations
measure \(D_M(z)/r_s\) and \(H(z)r_s\).
\end{exercise}

\begin{exercise}[Peculiar velocities in the local Hubble diagram]{Cambridge Part III Cosmology, examples}
At low redshift, include a source line-of-sight peculiar velocity in the
observed redshift and distance modulus.  Show why peculiar velocities dominate
the fractional distance error as \(z\to0\).
\end{exercise}
